\section{The comparison, frame by frame}
\label{sec:framewise}

Every figure so far reports a number pooled over a recording, and pooling is the
wrong summary for what this section measures. Two moving pictures and two still
figures are described below, all drawn from one recording that no policy was
trained on, and all showing the same instants. Figure~\ref{fig:framewise} gives
the contribution of touch through that recording, and
Figure~\ref{fig:gradcampeak} shows where each policy was looking at the moment
touch mattered most.

The procedure is the same for all four. The recording is replayed frame by
frame. At each frame every policy plans its next actions; then each input is
suppressed in turn and the policy plans again, and the distance between the two
plans is recorded as that input's contribution at that instant. The gradient
pictures are made the same way round but by a different route: instead of
suppressing an input, the sensitivity of the plan to each pixel of each camera
is computed directly and drawn over the image the policy was given. Because
every policy is asked about the same frame before any of them moves on, the
panels within one picture are one moment and not four.

That arrangement is enforced rather than hoped for. The comparison covers only
frames that every policy scored, so a policy that stopped early shortens the
comparison instead of shifting the others along it. Separate pictures played
side by side would drift apart within seconds, and a reader could no longer tell
whether two panels differ because the policies differ or because the moments do.

\begin{figure}[t]
  \centering
  \includegraphics[width=\textwidth]{framewise_shares.png}
  \caption{The share of the plan carried by the four fingertip cameras, through
  one held-back recording. Colour follows the family and a dashed line marks the
  cropped arm, so the pair within a colour is the crop comparison. Shares are
  normalised within a policy, so heights may be compared along a curve but not
  between colours.}
  \label{fig:framewise}
\end{figure}

\textbf{Touch swings far more within one recording than it differs between
policies.} For the action-chunking family its share runs from \num{0.009} to
\num{0.439} across ten seconds --- close to fiftyfold --- against a pooled mean
of \num{0.095}. The mean is not wrong, but it describes no instant that occurs:
the channel is very nearly ignored for the opening three seconds and carries
over two fifths of the plan at its peak. Any conclusion of the form ``the
tactile cameras contribute little'' drawn from a pooled figure is therefore a
statement about an average, not about the moments the sensor exists for.

\textbf{The two families do not read touch the same way.} The action-chunking
family reads it in bursts, its peak reaching four to five times its own mean;
the diffusion family reads it steadily, ranging eightfold where the other ranges
fiftyfold, and peaking at barely twice its mean. One family treats the
fingertips as an occasional check and the other as a continuous input, and that
difference is invisible in every pooled figure in this report. Why they differ
is not established here.

\begin{figure}[t]
  \centering
  \includegraphics[width=\textwidth]{gradcam_peak_touch.png}
  \caption{Where each policy looks at the instant touch carried the most, one
  row per policy and one column per camera. A cropped policy's tactile panels
  are drawn from the cropped image --- the image that policy actually receives
  --- which is why those rows appear magnified. Brightness is relative to each
  map's own mean, so it marks where a policy looks more than it looks
  everywhere, not how strongly it looks.}
  \label{fig:gradcampeak}
\end{figure}

Two things govern how Figure~\ref{fig:gradcampeak} should be read. A cropped
arm's tactile panels are drawn from the cropped image rather than the full
frame, because showing the full frame would illustrate attention to pixels that
policy never received. And the gradient maps differ in grid resolution between
the two families, finer for the action-chunking one, so a comparison down a
column is sound while a comparison across one confuses a policy's attention with
its map's coarseness. Within that limit the picture agrees with
Section~\ref{sec:cropping}: after cropping, bright regions sit noticeably closer
to the edges of the tactile panels.

These pictures cannot settle three things, and Section~\ref{sec:threats} carries
them. Suppressing an input measures how much a trained policy uses it, not
whether the task needs it. One recording is an illustration and not a sample.
And a policy can look in a sensible place and still plan badly --- where a
policy looks and whether it is right are different measurements, and only the
second is in Section~\ref{sec:results}.
